% ============================================================
% 第15章 分式 — 人教版八年级上册
% 编译命令：xelatex -interaction=nonstopmode 第15章_分式.tex
% ============================================================
\documentclass[12pt, a4paper]{ctexart}
\usepackage[margin=2.2cm]{geometry}
\usepackage{amsmath, amssymb}
\usepackage{enumitem}
\usepackage{tcolorbox}
\usepackage{booktabs}
\usepackage{array}
\usepackage{fancyhdr}
\usepackage{xcolor}
\usepackage{tikz}
\usepackage{hyperref}
\usepackage{graphicx}
\usepackage{multirow}
\usepackage{pgfplots}
\pgfplotsset{compat=1.18}

% 颜色定义
\definecolor{mainblue}{RGB}{41, 98, 168}
\definecolor{lightblue}{RGB}{230, 240, 255}
\definecolor{tipgreen}{RGB}{34, 139, 87}
\definecolor{lightgreen}{RGB}{235, 250, 240}
\definecolor{warnred}{RGB}{200, 50, 50}
\definecolor{lightred}{RGB}{255, 240, 240}
\definecolor{examorange}{RGB}{210, 120, 20}
\definecolor{lightorange}{RGB}{255, 248, 235}
\definecolor{apppurple}{RGB}{120, 60, 170}
\definecolor{lightpurple}{RGB}{245, 238, 255}

% 页眉页脚
\pagestyle{fancy}
\fancyhf{}
\fancyhead[L]{\small\textcolor{mainblue}{\textbf{人教版八年级上册 · 第15章\ 分式}}}
\fancyhead[R]{\small\textcolor{mainblue}{分式讲义}}
\fancyfoot[C]{\small —\ \thepage\ —}
\renewcommand{\headrulewidth}{0.6pt}
\renewcommand{\headrule}{\hbox to\headwidth{\color{mainblue}\leaders\hrule height \headrulewidth\hfill}}
\renewcommand{\footrulewidth}{0pt}

% tcolorbox 样式
\newtcolorbox{knowledgebox}[1][]{%
  colback=lightblue, colframe=mainblue,
  fonttitle=\bfseries\large, title=#1,
  arc=3mm, boxrule=0.8pt,
  left=8pt, right=8pt, top=6pt, bottom=6pt
}

\newtcolorbox{tipbox}[1][]{%
  colback=lightgreen, colframe=tipgreen,
  fonttitle=\bfseries, title=#1,
  arc=2mm, boxrule=0.6pt,
  left=6pt, right=6pt, top=4pt, bottom=4pt
}

\newtcolorbox{warnbox}[1][]{%
  colback=lightred, colframe=warnred,
  fonttitle=\bfseries, title=#1,
  arc=2mm, boxrule=0.6pt,
  left=6pt, right=6pt, top=4pt, bottom=4pt
}

\newtcolorbox{examplebox}[1][]{%
  colback=lightorange, colframe=examorange,
  fonttitle=\bfseries, title=#1,
  arc=3mm, boxrule=0.8pt,
  left=8pt, right=8pt, top=6pt, bottom=6pt
}

\newtcolorbox{appbox}[1][]{%
  colback=lightpurple, colframe=apppurple,
  fonttitle=\bfseries, title=#1,
  arc=3mm, boxrule=0.8pt,
  left=8pt, right=8pt, top=6pt, bottom=6pt
}

\setCJKmainfont{Hiragino Sans GB}
\setCJKsansfont{Heiti SC}

\begin{document}

% ============ 标题页 ============
\begin{center}
{\Huge\bfseries\textcolor{mainblue}{第15章\quad 分\quad 式}}\\[12pt]
{\Large\textcolor{gray}{讲\quad 义}}\\[8pt]
\rule{0.6\textwidth}{1pt}\\[8pt]
{\normalsize 适用：人教版\quad 八年级上册\quad |\quad 课时：4课时}
\end{center}
\vspace{10pt}

% ================================================================
% 第一部分：知识点讲解
% ================================================================
\section*{\textcolor{mainblue}{一、知识点讲解}}

% ---------- 15.1 分式的概念 ----------
\subsection*{15.1\quad 分式的概念}

\begin{knowledgebox}[定义：分式]
一般地，如果$A$、$B$表示两个整式，并且$B$中含有字母，那么式子$\dfrac{A}{B}$叫做\textbf{分式}。其中$A$叫做分子，$B$叫做分母。

\medskip
\textbf{关键要点：}
\begin{itemize}[leftmargin=2em]
  \item 分母$B$中必须含有字母（未知数），否则就是整式。
  \item 分式的本质仍然是除法运算，只不过被除数和除数用分数的形式表示。
  \item 分式$\dfrac{A}{B}$有意义的条件：$B \neq 0$（分母不为零）。
  \item 分式$\dfrac{A}{B}$无意义的条件：$B = 0$。
  \item 分式$\dfrac{A}{B} = 0$的条件：$A = 0$ 且 $B \neq 0$（分子为零且分母不为零）。
\end{itemize}
\end{knowledgebox}

\begin{warnbox}[易错提醒：分式的识别]
判断一个式子是否为分式，关键看\textbf{分母中是否含有字母}。

\medskip
例如：
\begin{itemize}[leftmargin=2em]
  \item $\dfrac{1}{x}$ 是分式（分母含字母$x$）。
  \item $\dfrac{x}{3}$ \textbf{不是}分式（分母是常数3，这是整式）。
  \item $\dfrac{x+1}{\pi}$ \textbf{不是}分式（$\pi$是常数，不是字母）。
  \item $\dfrac{x^2+1}{x-1}$ 是分式。
\end{itemize}
\end{warnbox}

\begin{tipbox}[判断分式"值为零"的三步法]
\begin{enumerate}[leftmargin=2em]
  \item 令分子$A=0$，解出字母的值；
  \item 将所得值代入分母$B$，检验$B \neq 0$；
  \item 若$B \neq 0$，则分式值为零；若$B = 0$，则分式无意义（舍去）。
\end{enumerate}
\end{tipbox}

% ---------- 15.2 分式的基本性质 ----------
\subsection*{15.2\quad 分式的基本性质}

\begin{knowledgebox}[分式的基本性质]
分式的分子与分母同时乘以（或除以）同一个不为零的整式，分式的值不变。

\[
\frac{A}{B} = \frac{A \cdot C}{B \cdot C} \quad (C \neq 0), \qquad \frac{A}{B} = \frac{A \div C}{B \div C} \quad (C \neq 0)
\]

\medskip
\textbf{应用一：约分}
\begin{itemize}[leftmargin=2em]
  \item 找分子、分母的\textbf{公因式}，然后约去。
  \item 公因式 = 系数的最大公约数 $\times$ 相同字母（取最低次幂）。
  \item 约分的结果必须是\textbf{最简分式}（分子与分母没有公因式，$1$除外）。
\end{itemize}

\medskip
\textbf{应用二：通分}
\begin{itemize}[leftmargin=2em]
  \item 找各分母的\textbf{最简公分母}，然后将每个分式化为同分母分式。
  \item 最简公分母 = 各分母系数的最小公倍数 $\times$ 所有出现的字母（取最高次幂）。
\end{itemize}
\end{knowledgebox}

\begin{warnbox}[易错提醒：约分时不能"跨项"约分]
约分只能约去分子和分母的\textbf{公因式}，不能只约分分子（或分母）中的一部分。

\medskip
\textbf{错误示例：} $\dfrac{x+2}{x+3} = \dfrac{2}{3}$\quad（错误！$x$不是分子和分母的公因式）

\medskip
\textbf{正确做法：}
\[
\frac{x^2 + 2x}{x^2 + 3x} = \frac{x(x+2)}{x(x+3)} = \frac{x+2}{x+3} \quad (x \neq 0)
\]
\end{warnbox}

% ---------- 15.3 分式的运算 ----------
\subsection*{15.3\quad 分式的运算}

\begin{knowledgebox}[分式的乘除法]
\textbf{乘法法则：}
\[
\frac{A}{B} \cdot \frac{C}{D} = \frac{A \cdot C}{B \cdot D}
\]
即分子乘以分子，分母乘以分母，然后约分化简。

\medskip
\textbf{除法法则：}
\[
\frac{A}{B} \div \frac{C}{D} = \frac{A}{B} \cdot \frac{D}{C} = \frac{A \cdot D}{B \cdot C} \quad (C \neq 0)
\]
除以一个分式等于乘以它的倒数。

\medskip
\textbf{乘方法则：}
\[
\left(\frac{A}{B}\right)^n = \frac{A^n}{B^n} \quad (B \neq 0,\; n\text{为正整数})
\]
\end{knowledgebox}

\begin{knowledgebox}[分式的加减法]
\textbf{同分母分式加减：}
\[
\frac{A}{C} \pm \frac{B}{C} = \frac{A \pm B}{C}
\]

\medskip
\textbf{异分母分式加减：}
\begin{enumerate}[leftmargin=2em]
  \item 找最简公分母；
  \item 将各分式通分；
  \item 按同分母分式加减法则计算；
  \item 结果化为最简分式。
\end{enumerate}
\end{knowledgebox}

\begin{tipbox}[分式运算的一般步骤]
\begin{enumerate}[leftmargin=2em]
  \item \textbf{先观察}：是否有可约分的项，能否先约分再计算。
  \item \textbf{乘除法}：先将除法统一变为乘法，再约分，最后相乘。
  \item \textbf{加减法}：先通分，再加减，最后化简。
  \item \textbf{结果要求}：化为最简分式或整式。
\end{enumerate}
\end{tipbox}

% ---------- 15.4 分式方程 ----------
\subsection*{15.4\quad 分式方程}

\begin{knowledgebox}[分式方程及其解法]
\textbf{定义：}分母中含有未知数的方程叫做\textbf{分式方程}。

\medskip
\textbf{解分式方程的一般步骤：}
\begin{enumerate}[leftmargin=2em]
  \item 找到各分母的\textbf{最简公分母}；
  \item 方程两边同时乘以最简公分母（\textbf{去分母}），将分式方程化为整式方程；
  \item 解这个整式方程；
  \item \textbf{验根}：将所得解代入最简公分母，若最简公分母不为零，则是原方程的解；若为零，则是\textbf{增根}，应舍去。
\end{enumerate}
\end{knowledgebox}

\begin{warnbox}[重要提醒：分式方程必须验根！]
去分母的过程相当于在方程两边同时乘以一个含有未知数的整式，这个整式有可能为零。因此解分式方程后\textbf{必须验根}。

\medskip
\textbf{验根方法：}将求得的解代入\textbf{最简公分母}，而不是代入原方程。若最简公分母为零，则该解为增根，应舍去。
\end{warnbox}

% ================================================================
% 第二部分：例题精讲
% ================================================================
\section*{\textcolor{examorange}{二、例题精讲}}

% --- 例题1 ---
\begin{examplebox}[【例题1】判断分式的类型及意义条件]
当$x$取何值时，分式$\dfrac{x-1}{x^2 - 4}$有意义？当$x$取何值时，分式的值为零？

\bigskip
\textbf{解：}

\textbf{(1) 有意义的条件：}

分母不为零，即 $x^2 - 4 \neq 0$。

解 $x^2 - 4 = 0$，得 $x = 2$ 或 $x = -2$。

$\therefore$ 当 $x \neq 2$ 且 $x \neq -2$ 时，分式有意义。

\medskip
\textbf{(2) 值为零的条件：}

分子为零且分母不为零。

令 $x - 1 = 0$，得 $x = 1$。

当 $x = 1$ 时，$x^2 - 4 = 1 - 4 = -3 \neq 0$。

$\therefore$ 当 $x = 1$ 时，分式的值为零。
\end{examplebox}

\begin{tipbox}[方法总结]
分式"有意义"只要求分母不为零；分式"值为零"要求分子为零\textbf{且}分母不为零。两个条件缺一不可，这是考试高频考点。
\end{tipbox}

% --- 例题2 ---
\begin{examplebox}[【例题2】分式的约分与化简]
化简下列各式：

\medskip
(1)\ $\dfrac{12a^3b^2}{-18a^2b^3}$

\medskip
(2)\ $\dfrac{x^2 - 9}{x^2 - 6x + 9}$

\medskip
\textbf{解：}

(1)\ $\dfrac{12a^3b^2}{-18a^2b^3} = \dfrac{12}{-18} \cdot \dfrac{a^3}{a^2} \cdot \dfrac{b^2}{b^3} = -\dfrac{2}{3} \cdot a \cdot \dfrac{1}{b} = -\dfrac{2a}{3b}$

\medskip
(2)\ $\dfrac{x^2 - 9}{x^2 - 6x + 9} = \dfrac{(x+3)(x-3)}{(x-3)^2} = \dfrac{x+3}{x-3} \quad (x \neq 3)$
\end{examplebox}

\begin{tipbox}[方法总结]
约分前先对分子、分母进行因式分解（提公因式、公式法等），找到公因式后再约去。注意约分后要标注使原分母为零的字母取值排除条件。
\end{tipbox}

% --- 例题3 ---
\begin{examplebox}[【例题3】分式的混合运算]
计算：

\medskip
(1)\ $\dfrac{2}{x-1} - \dfrac{1}{x+1}$

\medskip
(2)\ $\dfrac{x}{x-y} \cdot \dfrac{x^2-y^2}{x^2} \div \dfrac{x+y}{x}$

\bigskip
\textbf{解：}

(1)\ 最简公分母为 $(x-1)(x+1)$。
\begin{align*}
\frac{2}{x-1} - \frac{1}{x+1} &= \frac{2(x+1)}{(x-1)(x+1)} - \frac{1(x-1)}{(x-1)(x+1)} \\[6pt]
&= \frac{2(x+1) - (x-1)}{(x-1)(x+1)} \\[6pt]
&= \frac{2x + 2 - x + 1}{(x-1)(x+1)} \\[6pt]
&= \frac{x + 3}{(x-1)(x+1)}
\end{align*}

\medskip
(2)\ 先把除法变为乘法，再约分：
\begin{align*}
\frac{x}{x-y} \cdot \frac{x^2-y^2}{x^2} \div \frac{x+y}{x}
&= \frac{x}{x-y} \cdot \frac{(x+y)(x-y)}{x^2} \cdot \frac{x}{x+y} \\[6pt]
&= \frac{x \cdot (x+y)(x-y) \cdot x}{(x-y) \cdot x^2 \cdot (x+y)} \\[6pt]
&= 1
\end{align*}
\end{examplebox}

\begin{tipbox}[方法总结]
分式加减法的核心是\textbf{通分}，乘除法的核心是\textbf{约分}。混合运算中，先算乘除后算加减，能约分的先约分再相乘，可以大大简化计算。
\end{tipbox}

% --- 例题4 ---
\begin{examplebox}[【例题4】解分式方程]
解方程：$\dfrac{3}{x-2} + 1 = \dfrac{x-1}{2-x}$

\bigskip
\textbf{解：}

注意 $\dfrac{x-1}{2-x} = \dfrac{x-1}{-(x-2)} = -\dfrac{x-1}{x-2}$

原方程化为：
\[
\frac{3}{x-2} + 1 = -\frac{x-1}{x-2}
\]

两边同乘以 $(x-2)$（最简公分母）：
\[
3 + (x-2) = -(x-1)
\]

\[
3 + x - 2 = -x + 1
\]

\[
x + 1 = -x + 1
\]

\[
2x = 0 \implies x = 0
\]

\textbf{验根：}将 $x=0$ 代入最简公分母 $x-2 = 0-2 = -2 \neq 0$。

$\therefore$ $x=0$ 是原方程的解。
\end{examplebox}

\begin{tipbox}[方法总结]
解分式方程四步走：\textbf{找}公分母 $\to$ \textbf{去}分母 $\to$ \textbf{解}整式方程 $\to$ \textbf{验根}。注意当分母含 $a-x$ 形式时，先化为 $-(x-a)$，方便找公分母。
\end{tipbox}

% --- 例题5 ---
\begin{examplebox}[【例题5】分式方程有增根的情况]
关于$x$的方程 $\dfrac{x}{x-3} - 2 = \dfrac{m}{x-3}$ 有增根，求$m$的值。

\bigskip
\textbf{解：}

两边同乘以 $x - 3$：
\[
x - 2(x-3) = m
\]

\[
x - 2x + 6 = m
\]

\[
-x + 6 = m \implies x = 6 - m
\]

增根即使最简公分母为零的值，即 $x - 3 = 0$，$x = 3$。

令 $6 - m = 3$，解得 $m = 3$。

$\therefore$ 当 $m = 3$ 时，方程有增根 $x = 3$。
\end{examplebox}

\begin{tipbox}[方法总结]
"有增根"意味着解出的整式方程的根恰好使最简公分母为零。解题思路：先去分母得到含参数的整式方程，再令最简公分母为零求出增根，最后代入整式方程解出参数的值。
\end{tipbox}

% --- 例题6 ---
\begin{examplebox}[【例题6】分式化简求值]
先化简，再求值：$\left(\dfrac{x}{x-2} - \dfrac{x}{x+2}\right) \div \dfrac{4x}{x^2 - 4}$，其中 $x = -1$。

\bigskip
\textbf{解：}

先化简括号内的部分：
\[
\frac{x}{x-2} - \frac{x}{x+2} = \frac{x(x+2) - x(x-2)}{(x-2)(x+2)} = \frac{x^2+2x-x^2+2x}{(x-2)(x+2)} = \frac{4x}{(x-2)(x+2)}
\]

再除以 $\dfrac{4x}{x^2-4}$，即乘以倒数：
\[
\frac{4x}{(x-2)(x+2)} \div \frac{4x}{(x-2)(x+2)} = \frac{4x}{(x-2)(x+2)} \cdot \frac{(x-2)(x+2)}{4x} = 1
\]

将 $x = -1$ 代入化简结果 $1$，得原式 $= 1$。
\end{examplebox}

\begin{tipbox}[方法总结]
"先化简再求值"题型，关键是先把分式化为最简形式。化简时注意：(1)通分合并；(2)除法变乘法；(3)因式分解后约分。最后代入数值时要确保代入的是使原式有意义的值。
\end{tipbox}

% --- 例题7 ---
\begin{examplebox}[【例题7】综合提高——分式的条件求值]
已知 $\dfrac{1}{x} + \dfrac{1}{y} = 5$，求 $\dfrac{2x - 3xy + 2y}{x + 2xy + y}$ 的值。

\bigskip
\textbf{解：}

由 $\dfrac{1}{x} + \dfrac{1}{y} = 5$ 得 $\dfrac{y + x}{xy} = 5$，即 $x + y = 5xy$。

将 $x + y = 5xy$ 代入所求式：

\textbf{分子：}$2x - 3xy + 2y = 2(x+y) - 3xy = 2 \times 5xy - 3xy = 10xy - 3xy = 7xy$

\textbf{分母：}$x + 2xy + y = (x+y) + 2xy = 5xy + 2xy = 7xy$

$\therefore \dfrac{2x - 3xy + 2y}{x + 2xy + y} = \dfrac{7xy}{7xy} = 1$
\end{examplebox}

\begin{tipbox}[方法总结]
"条件求值"题型不要试图分别求出$x$和$y$的值，而应把条件整体代入。常见技巧：(1)从条件中提取$x+y$、$xy$等"整体"；(2)将所求式用相同的整体表示；(3)代入消元。
\end{tipbox}

% ================================================================
% 第三部分：应用题
% ================================================================
\section*{\textcolor{apppurple}{三、应用题}}

% --- 应用1 ---
\begin{appbox}[【应用1】工程问题]
一项工程，甲单独做需要$a$天完成，乙单独做需要$b$天完成。

(1) 甲、乙合做一天能完成全部工程的几分之几？

(2) 甲、乙合做需要多少天完成？

\bigskip
\textbf{解：}

(1) 甲每天完成$\dfrac{1}{a}$，乙每天完成$\dfrac{1}{b}$。

合做一天完成：$\dfrac{1}{a} + \dfrac{1}{b} = \dfrac{b+a}{ab}$

(2) 合做完成所需天数为：
\[
1 \div \frac{a+b}{ab} = \frac{ab}{a+b} \text{（天）}
\]

\textbf{答：}合做一天完成全部工程的$\dfrac{a+b}{ab}$，合做需要$\dfrac{ab}{a+b}$天完成。
\end{appbox}

% --- 应用2 ---
\begin{appbox}[【应用2】行程问题]
小明骑自行车从家到学校，去时速度为$v_1$千米/时，回来时速度为$v_2$千米/时（$v_1 \neq v_2$）。设家到学校的距离为$s$千米。

(1) 求小明往返的平均速度。

(2) 当 $v_1 = 12$，$v_2 = 8$ 时，求平均速度的值。

\bigskip
\textbf{解：}

(1) 去时用时 $\dfrac{s}{v_1}$ 小时，回来用时 $\dfrac{s}{v_2}$ 小时。

往返总路程 $= 2s$，总时间 $= \dfrac{s}{v_1} + \dfrac{s}{v_2} = \dfrac{s(v_1+v_2)}{v_1 v_2}$。

平均速度：
\[
\bar{v} = \frac{2s}{\dfrac{s(v_1+v_2)}{v_1 v_2}} = \frac{2v_1 v_2}{v_1 + v_2}
\]

(2) 当 $v_1 = 12$，$v_2 = 8$ 时：
\[
\bar{v} = \frac{2 \times 12 \times 8}{12 + 8} = \frac{192}{20} = 9.6 \text{（千米/时）}
\]

\textbf{答：}平均速度为 $\dfrac{2v_1 v_2}{v_1 + v_2}$ 千米/时，具体值为 $9.6$ 千米/时。
\end{appbox}

% --- 应用3 ---
\begin{appbox}[【应用3】浓度问题]
有含盐 $20\%$ 的盐水 $a$ 克，现在要配成含盐 $15\%$ 的盐水。

(1) 需要加多少克水？

(2) 当 $a = 100$ 时，需要加多少克水？

\bigskip
\textbf{解：}

(1) 设需要加水 $x$ 克。

盐的质量不变：$a \times 20\% = 0.2a$ 克。

加水后盐水总量 $= a + x$ 克，含盐比例为 $15\%$：
\[
\frac{0.2a}{a + x} = 0.15
\]

\[
0.2a = 0.15(a + x)
\]

\[
0.2a = 0.15a + 0.15x
\]

\[
0.05a = 0.15x
\]

\[
x = \frac{a}{3}
\]

(2) 当 $a = 100$ 时：$x = \dfrac{100}{3} \approx 33.3$（克）。

\textbf{答：}需要加水 $\dfrac{a}{3}$ 克，当 $a=100$ 时约需加 $33.3$ 克水。
\end{appbox}

% --- 应用4 ---
\begin{appbox}[【应用4】注水问题]
一个水池有甲、乙两个进水管。单独开甲管注满水池需要$a$小时，单独开乙管注满水池需要$b$小时（$a > b > 0$）。

(1) 两管同时开放，多少小时能注满水池？

(2) 若先开甲管 $2$ 小时，再两管同时开放，还需多少小时注满？

\bigskip
\textbf{解：}

(1) 甲管每小时注$\dfrac{1}{a}$，乙管每小时注$\dfrac{1}{b}$。

两管同时开放每小时注 $\dfrac{1}{a}+\dfrac{1}{b} = \dfrac{a+b}{ab}$。

注满水池所需时间：$1 \div \dfrac{a+b}{ab} = \dfrac{ab}{a+b}$（小时）。

(2) 先开甲管 $2$ 小时，注了 $\dfrac{2}{a}$。

剩余 $1 - \dfrac{2}{a} = \dfrac{a-2}{a}$。

两管同时开放注完剩余部分所需时间：
\[
\frac{a-2}{a} \div \frac{a+b}{ab} = \frac{a-2}{a} \cdot \frac{ab}{a+b} = \frac{b(a-2)}{a+b} \text{（小时）}
\]

\textbf{答：}两管同开$\dfrac{ab}{a+b}$小时注满；先开甲管$2$小时后再同开，还需$\dfrac{b(a-2)}{a+b}$小时。
\end{appbox}

% ================================================================
% 第四部分：课堂练习
% ================================================================
\section*{\textcolor{tipgreen}{四、课堂练习}}

\begin{enumerate}[leftmargin=2em]
  \item 当$x$为何值时，分式 $\dfrac{x+3}{x^2 - 1}$ 有意义？当$x$为何值时，分式的值为零？

  \item 化简：$\dfrac{m^2 - 4m + 4}{m^2 - 4}$

  \item 计算：$\dfrac{a}{a-b} - \dfrac{b}{a+b}$

  \item 解方程：$\dfrac{2}{x-1} = \dfrac{3}{x+1}$

  \item 先化简，再求值：$\left(1 + \dfrac{1}{x-1}\right) \div \dfrac{x}{x^2 - 1}$，其中 $x = 2$。
\end{enumerate}

\begin{knowledgebox}[参考答案]
\textbf{1.}\ 有意义的条件：$x^2 - 1 \neq 0$，即 $x \neq 1$ 且 $x \neq -1$。

值为零的条件：$x + 3 = 0$ 且 $x^2 - 1 \neq 0$，即 $x = -3$（代入验证：$x^2-1 = 8 \neq 0$，成立）。

\medskip
\textbf{2.}\ $\dfrac{m^2 - 4m + 4}{m^2 - 4} = \dfrac{(m-2)^2}{(m-2)(m+2)} = \dfrac{m-2}{m+2} \quad (m \neq 2)$

\medskip
\textbf{3.}\ $\dfrac{a}{a-b} - \dfrac{b}{a+b} = \dfrac{a(a+b) - b(a-b)}{(a-b)(a+b)} = \dfrac{a^2+ab-ab+b^2}{a^2-b^2} = \dfrac{a^2+b^2}{a^2-b^2}$

\medskip
\textbf{4.}\ 两边同乘 $(x-1)(x+1)$：$2(x+1) = 3(x-1)$，$2x+2=3x-3$，$x=5$。

验根：$(x-1)(x+1) = 4 \times 6 = 24 \neq 0$，$\therefore x = 5$。

\medskip
\textbf{5.}\ $\left(1 + \dfrac{1}{x-1}\right) \div \dfrac{x}{x^2-1} = \dfrac{x-1+1}{x-1} \cdot \dfrac{(x-1)(x+1)}{x} = \dfrac{x}{x-1} \cdot \dfrac{(x-1)(x+1)}{x} = x+1$

当 $x = 2$ 时，原式 $= 3$。
\end{knowledgebox}

\vfill
\begin{center}
\small\textcolor{gray}{本讲义仅供教学参考，例题可根据学生水平调整难度。}
\end{center}

\end{document}
